\section{Five episodes and their declared first stages}\label{sec:episodes}

Each episode enters the analysis as a declared first stage: published
estimates, imported with citations and attribution caveats, never
re-estimated. The five span the design's axes: negative and positive, price-led and
earnings-led, measured and projected, policy-made and
trade-transmitted. One is carried as a near-zero benchmark.
Figure~\ref{fig:episode-map} places the set on the paper's two axes and
Table~\ref{tab:first-stages} lists each declared first stage with its
source and epistemic status.

\begin{figure}[tbp]\centering
\includegraphics[width=\textwidth]{figures/fig_episode_map.png}
\caption{The episode set. Horizontal: the channel through which the
shock reaches households; vertical: cost or gain. Bubble area is
log-scaled in the gross shock. The 2022 energy episode (hatched, dashed
ring) is trade-\emph{transmitted} rather than a policy act; CPTPP (open,
dashed) is the near-zero benchmark.}
\label{fig:episode-map}
\end{figure}

% Generated by figures/make_figures.py -- do not edit by hand.
% Requires: \usepackage{booktabs}  (no other package needed)
% \rr = ragged-right inside a p{} cell, with \\ restored as the row
% terminator, so narrow columns neither overfull nor stretch-underfull.
\providecommand{\rr}{\raggedright\let\\\tabularnewline}

\begin{table}[tbp]
\centering
\footnotesize
\setlength{\tabcolsep}{3pt}
\renewcommand{\arraystretch}{1.15}
\caption{Declared first stages. Every price or earnings vector is an
imported published estimate; none is re-estimated in this paper.}
\label{tab:first-stages}
\begin{tabular}{@{}p{1.75cm} p{1.70cm} p{1.45cm} p{2.30cm} p{2.10cm} p{2.75cm} p{1.85cm}@{}}
\toprule
\rr Episode & \rr Date & \rr Type & \rr Channel & \rr First-stage source & \rr Key imported number & \rr Status \\
\midrule
\rr E1 TCA food NTBs & \rr Dec 2019--\newline Mar 2023 & \rr Policy & \rr Consumer prices (food) & \rr \citet{bakker2026} & \rr $+8\%$ on food prices (6\% low variant) & \rr Estimate\newline (ex post) \\[3pt]
\rr E2 Energy 2022--23 & \rr Oct 2022--\newline Mar 2023 & \rr Market & \rr Consumer prices (energy) & \rr \textsc{ofgem-cap}; \textsc{hmt-epg} & \rr Cap \pounds1{,}277 $\to$ \pounds3{,}549; EPG \pounds2{,}500 & \rr Observed\newline (statutory) \\[3pt]
\rr E3 US tariffs 2025 & \rr 2025 & \rr Policy\newline (foreign) & \rr Earnings only; no UK retaliation & \rr \citet{ahmadi2026} & \rr \pounds886m/yr gross earnings shock & \rr Estimate\newline (projection) \\[3pt]
\rr E4 India CETA & \rr 2026+ & \rr Policy & \rr Consumer prices (clothing, footwear, food) & \rr \citet{dbtindia2026} & \rr \pounds180m/yr duty cut on final goods & \rr Projection\newline (ex ante) \\[3pt]
\rr E5 CPTPP benchmark & \rr Long run & \rr Policy & \rr Aggregate GDP; no household mapping & \rr \citet{dbtcptpp2023} & \rr $+$\pounds2.0bn GDP ($+0.08\%$) & \rr Projection\newline (long run) \\[3pt]
\bottomrule
\end{tabular}

\vspace{2pt}
\begin{minipage}{\linewidth}\footnotesize
\emph{Notes.} Type distinguishes a policy instrument from a
trade-transmitted market price. Source entries are short cite keys; see the
bibliography. E3's consumer-price row is exactly zero (the UK imposed no
retaliatory tariffs), so its household incidence runs through earnings alone.
\end{minipage}
\end{table}


\subsection{Episode 1: the TCA border (2021--23), negative, price-led}

The Trade and Cooperation Agreement took effect in January 2021,
replacing single-market membership with a zero-tariff but
non-tariff-barrier-laden border. The price evidence is the cleanest of
any episode: exploiting cross-product variation in EU import intensity,
\citet{bakker2026} attribute roughly 8 percentage points of the 25 per
cent food price rise between December 2019 and March 2023 to Brexit
non-tariff barriers---a cumulative \pounds 250 per household on food
alone, \pounds 6.95 billion in aggregate (their household count;
\pounds 250 on the ONS count used here gives \pounds 7.1 billion), with pass-through of NTB costs
to consumers of 50--80 per cent and the largest effects in
high-EU-import products such as meat and dairy (around 10 points above
low-exposure comparators). The trade-volume record is equally stark:
relative UK imports from the EU fell about 25 per cent on entry
\citep{freeman2022}, and by 2024 goods exports were an estimated
\pounds 27 billion (6.4 per cent) below counterfactual, with some
16,400 small firms exiting EU exporting entirely
\citep{freeman2024}. The earnings side, by contrast, rests on
assumptions rather than estimates---the OBR's standing 15 per cent
trade-intensity and 4 per cent productivity wedges \citep{obrbrexit}---
so this episode's earnings channel is declared as an assumption-based
scenario and its price channel as an estimate-based one. Attribution
caveat: the price identification is relative (cross-product), so
economy-wide effects may be missed, and the counterfactual is contested
by critics of the OBR's trade assumptions.

\subsection{Episode 2: the terms-of-trade shock (2022--23), negative, price-led, very large}

The 2022 energy and import-price surge was, in the Bank of England's
description, a terms-of-trade shock of a size unseen since the
mid-1970s: the terms of trade deteriorated 6.2 per cent in the year to
2022Q3 \citep{onstot2023}, import prices rose 20 per cent relative to
output prices over two years \citep{broadbent2022}, the energy import
bill roughly doubled to \pounds 117 billion in 2022 \citep{imf2023}---the
\emph{increase} of roughly \pounds 60 billion being a transfer abroad
of about 2.5 per cent of GDP (derived arithmetic, labelled as such)---
and of the roughly 17 per cent CPI rise to 2022Q4, more than 70 per
cent (12.7 points) was imported \citep{dhingrapage2023}. The
household-level price vector is policy-mediated and well documented:
the typical annual energy bill's cap rose from about \pounds 1,277
(winter 2021--22) to \pounds 3,549 (October 2022), truncated to
\pounds 2,500 by the Energy Price Guarantee, at a discretionary cost
the OBR put at \pounds 27 billion for the EPG within \pounds 47
billion of price subsidies and \pounds 78 billion of total energy
support \citep{obrenergy2022}. Counterfactual no-policy welfare losses
average 6 per cent of after-tax income, reaching 11 per cent at the
95th percentile of exposure \citep{levell2025}, against energy budget
shares of roughly 11 per cent in the second income decile versus 6 per
cent in the ninth \citep{rf2022}. This is the strongest first stage of
the five---multiple independent official estimates of both the
aggregate shock and the retail price vector, plus a published
imported/domestic decomposition---and the design decision it forces is
constitutive for this paper: the shock is declared \emph{gross}, and
the Energy Price Guarantee is modelled as what it was, a discretionary
rulebook intervention, so that automatic and discretionary cushioning
can be separated.

\subsection{Episode 3: the US tariffs (2025--26), negative, earnings-led}

The April 2025 US schedule (10 per cent baseline, 25 per cent autos and
steel) and its partial mitigation are the subject of the companion
study \citep{ahmadi2026}, which supplies this episode's earnings first
stage and its second-stage results: a tariff-calibrated earnings loss of
roughly \pounds 886 million a year, cushioned at 43.9 per cent when
delivered as broad wage cuts and 36.6 per cent as displacement. The
external anchors are the OBR's March 2025 scenarios (import-demand
elasticity 0.4; peak GDP effects of $-0.6$ to $-1.0$ per cent across
scenarios) \citep{obr2025} and the Bank of England's effective-tariff
estimates (11 per cent at April 2025, about 9 per cent after the
Economic Prosperity Deal) \citep{boempr2025}. The UK consumer side of
this episode is, to date, literally zero: the UK has not retaliated, so
no tariff-induced consumer price vector exists---an empirical fact the
comparative design records rather than an omission.

\subsection{Episode 4: the India agreement (2026--), positive, both channels}

The UK--India Comprehensive Economic and Trade Agreement entered into
force on 15 July 2026. Its first stage is declared from the
government's own impact assessment \citep{dbtindia2026}: long-run GDP
$+\pounds 4.8$ billion a year ($+0.13$ per cent by 2040), real wages
$+0.19$ per cent ($\approx\pounds 2.2$ billion), and---on a different
clock, at entry into force---consumer duty reductions of \pounds 220
million a year, \pounds 180 million of it on final goods, concentrated
in clothing (\pounds 132 million), footwear (\pounds 13 million) and
food and beverages (\pounds 12 million). Sectoral GVA effects give the
earnings channel its geography: machinery and equipment $+1.65$ per
cent, textiles and apparel $-0.68$ per cent. Two dials are declared
rather than estimated, following the IA's own caveats: preference
utilisation (the duty figures assume 100 per cent on static 2022
trade) and consumer pass-through of the duty savings. The honest
framing---and the reason a government projection can anchor a
conditional design---is that the exercise computes what the
government's own central estimates imply for households, which its
impact assessment does not report and its regulatory scrutineer asked
for \citep{rpc2026}.

\subsection{Episode 5: CPTPP accession (2024--), the near-zero benchmark}

The UK's accession to the CPTPP (in force December 2024) carries a
central government estimate of $+\pounds 2.0$ billion a year in the
long run---0.08 per cent of GDP, about \pounds 70 per household per
year before any incidence structure---because the UK already had
agreements with nine of the eleven members \citep{dbtcptpp2023}. It is
included precisely because it is nearly nothing: a data-limitation
benchmark (the replication artifacts label it accordingly) whose honest treatment---an aggregate divided by a household
count, and no more---marks the boundary of what a declared first stage
can support.

\subsection{What the five episodes jointly reveal about the record}

One pattern runs across the dossier and should be stated carefully,
because its naive version is mechanical: past shocks naturally have
ex-post estimates and new agreements ex-ante projections. The
non-mechanical fact is on the evaluation side---no UK trade agreement
has ever been ex-post evaluated at the household level, although CPTPP
has now been in force long enough for the exercise, and the India
assessment's own regulatory scrutineer flagged the consumer analysis
as underdeveloped \citep{rpc2026}. The comparative design keeps the
epistemic split explicit throughout---estimate-based and
projection-based first stages are labelled in every table---and the
episode set is typed policy/market, with the terms-of-trade episode
carried as a trade-transmitted market shock, never as a policy act.
